\documentclass[UTF8,openany,oneside,zihao=-4]{ctexbook}

\usepackage[a4paper,margin=2cm,headheight=14pt]{geometry}
\usepackage{amsmath,amssymb}
\usepackage{xcolor}
\usepackage{enumitem}
\usepackage{fancyhdr}
\usepackage{listings}
\usepackage{titlesec}
\usepackage{tcolorbox}
\tcbuselibrary{skins,breakable}
\usepackage{longtable}
\usepackage{booktabs}
\usepackage{array}
\usepackage{calc}
\usepackage[
  hidelinks,
  unicode,
  bookmarksopen=true,
  bookmarksnumbered=true,
  pdfpagelabels=true
]{hyperref}

\providecommand{\tightlist}{%
  \setlength{\itemsep}{0pt}\setlength{\parskip}{0pt}}
\providecommand{\passthrough}[1]{#1}
\providecommand{\real}[1]{#1}
\newfontfamily\newboardcodefont{Cascadia Mono}
\newcommand{\inputnewboard}[1]{%
  {%
    \lstset{basicstyle=\newboardcodefont\small}%
    \input{#1}%
  }%
}

% ========== 颜色 ==========
\definecolor{codeblue}{HTML}{0550AE}
\definecolor{codegray}{HTML}{6E7781}
\definecolor{codebg}{HTML}{F6F8FA}
\definecolor{warnred}{HTML}{CF222E}
\definecolor{fixgreen}{HTML}{116329}
\definecolor{teal}{HTML}{0E6E6E}

% ========== 页眉页脚 ==========
\pagestyle{fancy}
\fancyhf{}
\fancyhead[L]{\small\textbf{算法板子}}
\fancyhead[R]{\small\thepage}
\renewcommand{\headrulewidth}{0.3pt}

% ========== 段落与列表 ==========
\setlength{\parindent}{0pt}
\setlength{\parskip}{0.4em}
\setlist[itemize]{leftmargin=1.8em,itemsep=0.15em,topsep=0.15em}
\setlist[enumerate]{leftmargin=2em,itemsep=0.2em,topsep=0.2em}

% ========== 标题样式 ==========
\titleformat{\chapter}{\LARGE\bfseries}{\thechapter}{0.7em}{}[\vspace{0.2em}\hrule]
\titleformat{\section}{\large\bfseries\color{codeblue}}{\thesection}{0.5em}{}
\titleformat{\subsection}{\normalsize\bfseries\color{teal}}{\thesubsection}{0.5em}{}

% ========== 代码样式 ==========
\lstdefinestyle{cpp}{
  language=C++,
  basicstyle=\ttfamily\small,
  keywordstyle=\color{codeblue}\bfseries,
  commentstyle=\color{codegray},
  stringstyle=\color{warnred},
  backgroundcolor=\color{codebg},
  frame=single,
  rulecolor=\color{codebg},
  breaklines=true,
  columns=fullflexible,
  keepspaces=true,
  showstringspaces=false,
  tabsize=4,
  literate=
    {≤}{{$\leq$}}1 {≥}{{$\geq$}}1 {≡}{{$\equiv$}}1
    {∤}{{$\nmid$}}1 {φ}{{$\varphi$}}1 {∏}{{$\prod$}}1
    {μ}{{$\mu$}}1 {α}{{$\alpha$}}1 {≈}{{$\approx$}}1 {∨}{{$\lor$}}1
    {├}{{|}}1 {└}{{`}}1 {─}{{--}}1
}
\lstset{style=cpp}

% ========== 提示框 ==========
\newtcolorbox{warn}{
  colback=red!3,colframe=warnred,
  fonttitle=\bfseries,title={\color{warnred}注意},
  boxrule=0.4pt,left=4pt,right=4pt,top=2pt,bottom=2pt,
  breakable
}
\newtcolorbox{tip}{
  colback=green!3,colframe=fixgreen,
  fonttitle=\bfseries,title={\color{fixgreen}用法/扩展},
  boxrule=0.4pt,left=4pt,right=4pt,top=2pt,bottom=2pt,
  breakable
}
\newtcolorbox{note}{
  colback=blue!3,colframe=codeblue,
  boxrule=0.4pt,left=4pt,right=4pt,top=2pt,bottom=2pt,
  breakable
}

\begin{document}

% ========== 封面 ==========
\begin{titlepage}
\thispdfpagelabel{封面}
\centering
\vspace*{3cm}
{\Huge\bfseries 算法板子\par}
\vspace{0.6cm}
{\Large 竞赛代码模板完整大版本\par}
\vspace{1.5cm}
{\normalsize 完整覆盖 · 专题详解 · 统一目录\par}
\vfill
{\large \today\par}
\end{titlepage}
\clearpage

\tableofcontents
\clearpage

\inputnewboard{remake/large/00_总览.tex}

% ==========================================================
\chapter{竞赛模板}
% ==========================================================

\begin{warn}
\begin{itemize}
  \item \texttt{\#define int long long} 之后 main 必须写 \texttt{signed main}，否则编译错误。
  \item 位运算中常数 \texttt{1} 默认是 32 位，当 $i \ge 31$ 时 \texttt{1 << i} 溢出变负数，必须写 \texttt{1LL << i}。
  \item \texttt{inf = 2e18} 做加法时可能溢出（如 \texttt{inf + x} 爆 long long），松弛前先判 \texttt{if(dist[u] != inf)}。
\end{itemize}
\end{warn}

\begin{lstlisting}
#include <bits/stdc++.h>
using namespace std;
#define int long long
#define endl '\n'

// 变量：inf=表示不可达状态的正无穷哨兵值 inf；参与加法前先判断当前值不是 inf。
const int inf = 2e18;
// 变量：N=当前代码允许使用的最大下标上界 N。
const int N = 1e6 + 10;

// 接口：solve()：读取一个测试用例的规模 n，并作为题目核心逻辑的单测入口；无返回值。
void solve() {
    // 变量：n=当前测试用例读入的元素个数 n。
    int n;
    cin >> n;
}

// 接口：main()：开启快速 I/O，读入测试数 t，并为每个测试调用一次 solve；入口函数；每组测试前按注释清空全局状态。
signed main() {
    ios::sync_with_stdio(0);
    cin.tie(0);
    // 变量：t=当前程序要执行的测试用例数量 t。
    int t = 1;
    cin >> t;
    while (t--) {
        solve();
    }
    return 0;
}
\end{lstlisting}

\begin{tip}
\begin{itemize}
  \item 不需要多测时把 \lstinline|cin >> t| 注释掉即可。
  \item 多测时注意：solve() 里不要中途 \lstinline|return|（会导致后续输入串流）；全局数组/容器必须清空干净。
  \item 局部变量必须初始化！\lstinline|int sum;| 在 OJ 上可能是乱码，必须写 \lstinline|int sum = 0;|。
\end{itemize}
\end{tip}

\inputnewboard{remake/large/01_模板.tex}

% ==========================================================
\chapter{STL 常用函数速查}
% ==========================================================

\section{容器}

\subsection{vector}
\begin{lstlisting}
vector<int> a(n, 0);       // n 个元素初始化为 0
a.push_back(x);            // 尾部插入
a.pop_back();              // 尾部删除
a.size();                  // 元素个数（返回 size_t，比较时注意）
a.empty();                 // 是否为空
a.clear();                 // 清空
a.resize(m, val);          // 重设大小
a.erase(a.begin() + i);   // 删除第 i 个元素 O(n)
// 去重：
sort(a.begin(), a.end());
a.erase(unique(a.begin(), a.end()), a.end());
\end{lstlisting}

\subsection{array}
\begin{lstlisting}
array<int, N> a;           // 固定大小数组，N 必须是编译期常量
array<int, 5> a = {1, 2, 3, 4, 5};  // 初始化
array<int, 5> a = {};      // 全部初始化为 0
a.size();                  // 元素个数（编译期确定，始终为 N）
a.fill(val);               // 全部填充为 val
a.front();  a.back();      // 首尾元素
a.data();                  // 底层数组指针
a[i];                      // 随机访问 O(1)
// 支持比较运算符（按字典序），可以直接当 map 的 key
// 与 C 数组的区别：可以赋值、比较、作为函数返回值
\end{lstlisting}

\begin{warn}
\begin{itemize}
  \item \lstinline|array| 不能用变量做大小，\lstinline|array<int, n>| 编译错误（n 必须是 \lstinline|constexpr|）。
  \item 局部 \lstinline|array| 不会自动初始化为 0！必须写 \lstinline|array<int, 5> a = {};| 或 \lstinline|a.fill(0);|。
  \item 竞赛中常用场景：需要把固定长度的数组当 \lstinline|map| / \lstinline|set| 的 key 时，用 \lstinline|array| 比手写 \lstinline|struct| 方便。
\end{itemize}
\end{warn}

\subsection{string}
\begin{lstlisting}
// 变量：s=从标准输入读取、供当前字符串算法处理的字符串 s。
string s;
getline(cin, s);           // 读一整行（注意前面 cin>>n 后要 cin.ignore()）
s.substr(pos, len);        // 从 pos 开始取 len 个字符
s.find("abc");             // 查找子串，找不到返回 string::npos
s += "xyz";                // 拼接
to_string(123);            // int -> string
stoi("123");               // string -> int
stoll("123456789012");     // string -> long long
\end{lstlisting}

\subsection{set / multiset}
\begin{lstlisting}
// 变量：s=用于演示有序集合接口的 set 容器 s。
set<int> s;
s.insert(x);              // 插入
s.erase(x);               // 删除所有值为 x 的元素
s.count(x);               // 0 或 1
s.lower_bound(x);         // >= x 的第一个迭代器
s.upper_bound(x);         // > x 的第一个迭代器
// multiset 删除单个元素：
// 变量：it=容器 ms 中当前查找结果的迭代器 it。
auto it = ms.find(x);
if (it != ms.end()) ms.erase(it); // 只删一个！
// 如果写 ms.erase(x) 会删掉所有等于 x 的元素！
\end{lstlisting}

\subsection{map}
\begin{lstlisting}
// 变量：mp=把 int 键映射到 int 值并保持键有序的 map 容器 mp。
map<int, int> mp;
mp[key] = val;             // 插入/修改
mp.count(key);             // 是否存在
mp.erase(key);             // 删除
// 遍历：
for (auto& [k, v] : mp) { ... }
// 注意：mp[key] 如果 key 不存在会自动插入默认值 0！
// 只查询用 mp.count(key) 或 mp.find(key)
\end{lstlisting}

\subsection{unordered\_map / unordered\_set}
\begin{lstlisting}
unordered_map<int, int> mp;  // 哈希表，平均 O(1) 查找/插入/删除
mp[key] = val;              // 插入/修改（同 map）
mp.count(key);              // 是否存在
mp.erase(key);              // 删除
for (auto& [k, v] : mp) { ... }  // 遍历（无序！）

// 变量：us=当前无序集合 us。
unordered_set<int> us;
us.insert(x);
us.count(x);                // 0 或 1
us.erase(x);

// 与 map/set 的区别：
// 1. 无序，不支持 lower_bound / upper_bound
// 2. 平均 O(1)，但最坏 O(n)（哈希冲突）
// 3. 不能直接用 pair/vector 做 key（需要自定义哈希）
\end{lstlisting}

\begin{warn}
\begin{itemize}
  \item \textbf{被卡哈希}：Codeforces 等平台可以构造数据卡 \lstinline|unordered_map| 的默认哈希，导致退化为 $O(n^2)$。
  \item 防卡方法：自定义哈希函数加随机种子：
\end{itemize}
\end{warn}

\begin{lstlisting}
struct custom_hash {
    // 接口：operator()：返回键值的 size_t 哈希值，供 unordered_map/unordered_set 去重；返回类型 size_t。
    size_t operator()(int x) const {
        static const size_t FIXED_RANDOM =
            chrono::steady_clock::now().time_since_epoch().count();
        x ^= FIXED_RANDOM;
        return x ^ (x >> 16);
    }
};
// 变量：mp=使用 custom_hash 把 int 键映射到 int 值的哈希表 mp。
unordered_map<int, int, custom_hash> mp;
\end{lstlisting}

\begin{tip}
\begin{itemize}
  \item 用 \lstinline|mp.reserve(n)| 预分配空间、\lstinline|mp.max_load_factor(0.25)| 降低负载因子，可以显著提速。
  \item 需要有序遍历或 \lstinline|lower_bound| 时，必须用 \lstinline|map| / \lstinline|set|。
  \item \lstinline|unordered_map| 的 \lstinline|[]| 操作同样会自动插入默认值，注意与 \lstinline|map| 一样的坑。
\end{itemize}
\end{tip}

\subsection{priority\_queue}
\begin{lstlisting}
// 大根堆（默认）
// 变量：pq=当前优先队列 pq；堆顶含义见比较器。
priority_queue<int> pq;
// 小根堆
// 变量：pq=当前优先队列 pq；堆顶含义见比较器。
priority_queue<int, vector<int>, greater<int>> pq;
pq.push(x);
pq.top();    // 堆顶（不弹出）
pq.pop();    // 弹出堆顶
\end{lstlisting}

\subsection{deque}
\begin{lstlisting}
// 变量：dq=当前双端队列 dq。
deque<int> dq;
dq.push_front(x); dq.push_back(x);
dq.pop_front();   dq.pop_back();
dq.front();       dq.back();
dq[i];            // 支持随机访问
\end{lstlisting}

\subsection{stack}
\begin{lstlisting}
// 变量：stk=当前栈 stk。
stack<int> stk;
stk.push(x);      // 压栈
stk.pop();         // 弹栈（无返回值！）
stk.top();         // 取栈顶元素
stk.empty();       // 是否为空
stk.size();        // 元素个数
// 注意：stack 没有 clear()，清空方式：
// while (!stk.empty()) stk.pop();
// 或直接 stk = stack<int>();
\end{lstlisting}

\subsection{list}
\begin{lstlisting}
list<int> lst;
lst.push_back(x);  lst.push_front(x);   // 头尾插入 O(1)
lst.pop_back();    lst.pop_front();      // 头尾删除 O(1)
lst.front();       lst.back();           // 访问头尾
lst.insert(it, x);                       // 在迭代器 it 前插入 O(1)
lst.erase(it);                           // 删除迭代器指向的元素 O(1)
lst.size();        lst.empty();
lst.remove(val);                         // 删除所有值为 val 的元素
lst.sort();                              // 链表专用排序（归并）O(n log n)
lst.unique();                            // 去除相邻重复元素（需先排序）
lst.merge(lst2);                         // 合并两个有序链表
lst.reverse();                           // 反转链表
lst.splice(it, lst2);                    // 将 lst2 全部移接到 it 前面 O(1)
// 注意：list 不支持随机访问（没有 []），遍历只能用迭代器
// 适用场景：频繁在中间插入/删除，且不需要随机访问
\end{lstlisting}

\section{常用算法}

\begin{lstlisting}
// 排序
sort(a.begin(), a.end());                    // 升序
sort(a.begin(), a.end(), greater<int>());    // 降序

// 二分查找（要求有序）
lower_bound(a.begin(), a.end(), x);  // >= x 的第一个位置
upper_bound(a.begin(), a.end(), x);  // > x 的第一个位置

// 去重
unique(a.begin(), a.end());  // 返回去重后的尾迭代器

// 全排列
next_permutation(a.begin(), a.end()); // 下一个排列，到最后返回 false

// 最值
*max_element(a.begin(), a.end());
*min_element(a.begin(), a.end());

// 累加
accumulate(a.begin(), a.end(), 0LL); // 注意初始值类型！
\end{lstlisting}

\begin{warn}
\begin{itemize}
  \item 自定义 \lstinline|sort| 的比较函数必须用严格小于 \lstinline|<|，不能用 \lstinline|<=|，否则破坏严格弱序，导致随机 RE/WA。
  \item \lstinline|accumulate| 初始值写 \lstinline|0| 会按 int 累加溢出，必须写 \lstinline|0LL|。
  \item \lstinline|vector.size()| 返回无符号数，\lstinline|a.size() - 1| 当 size 为 0 时会下溢变成极大值！
\end{itemize}
\end{warn}

\section{pair / tuple}

\begin{lstlisting}
// pair
// 变量：p=用于演示 pair 接口的整数对 p。
pair<int, int> p = {1, 2};
auto [x, y] = p;  // C++17 结构化绑定
p.first; p.second; // 访问元素
// pair 支持比较运算符（先比 first，再比 second）
// 可以直接作为 map 的 key、放入 set、sort

// make_pair
auto p2 = make_pair(3, 4); // 自动推导类型

// tuple（三元及以上）
tuple<int, int, int> t = {1, 2, 3};
auto [a, b, c] = t;  // C++17 结构化绑定
get<0>(t); get<1>(t); get<2>(t); // 访问元素
// 变量：t2=用于演示 tuple 构造与字典序比较的三元组 t2。
auto t2 = make_tuple(1, 2, 3);
// tuple 也支持比较运算符（按字典序）
// 竞赛中三元排序直接用 tuple 最方便
\end{lstlisting}

\begin{tip}
\begin{itemize}
  \item 多关键字排序用 \lstinline|vector<tuple<int,int,int>>| 比写结构体更快。
  \item \lstinline|tie(a, b, c) = make_tuple(x, y, z)| 可以一行赋值多个变量。
  \item C++17 的结构化绑定 \lstinline|auto [a, b, c] = t| 最简洁，推荐使用。
\end{itemize}
\end{tip}

\section{bitset}

\begin{note}
\lstinline|bitset<N>| 是固定大小的位集合，N 必须是编译期常量。空间 $N/8$ 字节，位运算操作是 $O(N/64)$（64 位压缩）。
\end{note}

\begin{lstlisting}
bitset<1000> bs;          // 全部初始化为 0
bitset<8> bs2("10101010"); // 从字符串初始化（注意：字符串最右对应第 0 位）
bitset<8> bs3(0xAA);      // 从整数初始化

// 单位操作
bs.set(i);        // 第 i 位置 1
bs.set(i, 0);     // 第 i 位置 0
bs.reset(i);      // 第 i 位置 0
bs.flip(i);       // 第 i 位取反
bs.test(i);       // 返回第 i 位的值（bool）
bs[i];            // 同 test(i)，但不做越界检查

// 全局操作
bs.set();         // 全部置 1
bs.reset();       // 全部置 0
bs.flip();        // 全部取反
bs.count();       // 1 的个数（popcount）
bs.size();        // 总位数 N
bs.any();         // 是否存在至少一个 1
bs.none();        // 是否全为 0
bs.all();         // 是否全为 1（C++11）

// 位运算（O(N/64)，比手写快很多）
bitset<1000> a, b;
a & b;  a | b;  a ^ b;  ~a;   // 返回新 bitset
a &= b; a |= b; a ^= b;       // 原地修改
a <<= k; a >>= k;             // 左移/右移 k 位
(a & b).count();               // 两个集合的交集大小

// 转换
bs.to_ulong();    // 转为 unsigned long（N <= 32 时）
bs.to_ullong();   // 转为 unsigned long long（N <= 64 时）
bs.to_string();   // 转为 "010101..." 字符串

// 查找（非标准但 GCC 支持）
bs._Find_first();  // 第一个 1 的位置（没有则返回 N）
bs._Find_next(i);  // i 之后第一个 1 的位置（没有则返回 N）
// 枚举所有 1 的位置：
// 变量：i=动态 bitset 中当前找到的置位下标 i。
for (int i = bs._Find_first(); i < (int)bs.size(); i = bs._Find_next(i)) {
    // 处理位置 i
}
\end{lstlisting}

\begin{warn}
\begin{itemize}
  \item \lstinline|bitset| 的大小 N 必须是编译期常量，不能用变量！
  \item \lstinline|bs[i]| 的 $i$ 从 0 开始，第 0 位是最低位。
  \item 从字符串构造时，字符串的最后一个字符对应第 0 位（最低位），容易搞反。
  \item \lstinline|_Find_first| 和 \lstinline|_Find_next| 是 GCC 扩展，非标准，但竞赛中可用。
\end{itemize}
\end{warn}

\begin{tip}
\begin{itemize}
  \item \textbf{bitset 优化 DP}：将布尔 DP 数组改用 bitset，转移用位移 + 或运算，复杂度除以 64。经典例题：背包可行性判断 \lstinline|bs |= (bs << w[i])|。
  \item \textbf{bitset 求交集/汉明距离}：\lstinline|(a ^ b).count()| 即汉明距离，\lstinline|(a & b).count()| 即共同元素个数。
  \item \textbf{bitset 优化图算法}：传递闭包 $O(n^3/64)$，二部图匹配加速等。
  \item \textbf{bitset 加速暴力}：$n \le 40000$ 的 $O(n^2)$ 布尔运算用 bitset 可以压到 $O(n^2/64)$，约等于 $n \le 5000$ 的常规暴力速度。
\end{itemize}
\end{tip}

\subsection{bitset 优化背包 / 子集和可行性}

\begin{note}
判断若干物品能凑出哪些重量（布尔 DP），把 \lstinline|dp[j]| 改为 bitset 的第 $j$ 位，转移用左移 + 或，复杂度 $O(nV/64)$。
\end{note}

\begin{lstlisting}
// 数值域：0<=w[i]<MAXV；左移超出 MAXV 的位按 bitset 规则丢弃。
// bs[j] = 1 表示重量 j 可以被凑出
bitset<MAXV> bs;
bs[0] = 1;                 // 重量 0 总能凑出（什么都不选）
// 变量：i=bitset 0/1 背包当前加入的物品下标 i。
for (int i = 0; i < n; i++)
    bs |= bs << w[i];      // 01 背包：每个物品做一次整体左移再或
// bs[j] 为 1 即容量 j 可达；bs.count() 为可达重量种类数

// 完全背包可行性：物品可重复选，按位逐位递推
// for (int i = 0; i < n; i++)
//     for (int j = w[i]; j < MAXV; j++) if (bs[j-w[i]]) bs[j] = 1;
\end{lstlisting}

\begin{warn}
\begin{itemize}
  \item \lstinline|MAXV| 必须 $\ge$ 最大可能重量 $+1$，且是编译期常量。
  \item \lstinline|bs |= bs << w[i]| 一次只把每个物品计入一次，是 01 背包语义；完全背包需另写逐位递推。
\end{itemize}
\end{warn}

\subsection{bitset 优化传递闭包（Floyd）}

\begin{note}
求有向图可达性矩阵，$O(n^3)$ 的 Floyd 用 bitset 整行或运算压到 $O(n^3/64)$。
\end{note}

\begin{lstlisting}
// 接口：transitive_closure(reach, n)：原地计算前 n 个顶点的传递闭包。
// 参数：reach 初始含自环和原图边；n 为实际顶点数且 n<=MAXN。
template<size_t MAXN>
void transitive_closure(array<bitset<MAXN>, MAXN>& reach, int n) {
    // 变量：k=当前允许作为中转点的顶点 k。
    for (int k = 0; k < n; k++)
        // 变量：i=当前起点 i。
        for (int i = 0; i < n; i++)
            if (reach[i][k])
                reach[i] |= reach[k];
}
\end{lstlisting}

\subsection{bitset 优化字符串匹配（Shift-And）}

\begin{note}
在文本 $t$ 中查找模式串 $p$（长 $m$），用 bitset 维护"当前已匹配前缀"集合，$O(nm/64)$。
\end{note}

\begin{lstlisting}
// B[c] 第 j 位为 1 表示 p[j]==c
bitset<MAXM> B[256], D;
// 变量：matches=模式串在文本中的全部 0-based 起始位置。
vector<int> matches;
// 变量：j=Shift-And 模式串 p 当前写入位掩码的字符位置 j。
for (int j = 0; j < m; j++) B[(unsigned char)p[j]][j] = 1;
// 变量：i=Shift-And 扫描主串 t 的当前字符位置 i。
for (int i = 0; i < (int)t.size(); i++) {
    D = ((D << 1) | bitset<MAXM>(1)) & B[(unsigned char)t[i]];
    if (D[m - 1]) {
        matches.push_back(i - m + 1);
    }
}
\end{lstlisting}

\begin{tip}
Shift-And 的状态 \lstinline|D[j]=1| 表示 \lstinline|p[0..j]| 是 \lstinline|t| 当前结尾的后缀。把 \lstinline{|} 换成"或非"、判 \lstinline|D[m-1]==0| 即为 Shift-Or 变体，原理相同。
\end{tip}

\section{\_\_int128}

\begin{note}
GCC 扩展的 128 位整数，范围约 $[-1.7 \times 10^{38}, 1.7 \times 10^{38}]$。用于中间计算结果超过 long long 但最终结果不超过的情况。
\end{note}

\begin{lstlisting}
// 变量：a=演示 __int128 算术与模乘的 128 位整数 a。
__int128 a = 1;
// 不能直接 cin/cout，需要手写 IO
// 接口：print128(x)：十进制输出一个 __int128；不返回值。
// 参数：待十进制输出的 __int128 整数 x。
void print128(__int128 x) {
    if (x < 0) { putchar('-'); x = -x; }
    if (x > 9) print128(x / 10);
    putchar(x % 10 + '0');
}

// 接口：read128()：从标准输入读取并返回一个 __int128；返回类型 __int128。
__int128 read128() {
    // 变量：x=十进制读入过程中累计绝对值的 __int128 结果 x；f=读入 __int128 时记录正负号的符号 f；c=当前从输入读取的字符 c。
    __int128 x = 0; int f = 1; char c = getchar();
    while (c < '0' || c > '9') { if (c == '-') f = -1; c = getchar(); }
    while (c >= '0' && c <= '9') { x = x * 10 + c - '0'; c = getchar(); }
    return x * f;
}

// 典型应用：大数乘法取模，避免 long long 乘法溢出
// 接口：mul_mod(a, b, mod)：计算 a*b mod mod，并避免中间乘法溢出；返回类型 int。
// 参数：模乘的第一个整数因子 a；模乘的第二个整数因子 b；正模数 mod。
int mul_mod(int a, int b, int mod) {
    return (__int128)a * b % mod;
}
\end{lstlisting}

\begin{warn}
\begin{itemize}
  \item \lstinline|__int128| 不能直接用 \lstinline|cin/cout/printf|，必须手写 IO。
  \item \lstinline|__int128| 是 GCC 扩展，MSVC 不支持。竞赛 OJ 一般都支持。
  \item 不能用于 \lstinline|bitset| 的下标或 \lstinline|vector| 的大小。
\end{itemize}
\end{warn}

\section{\_\_builtin 系列}

\begin{lstlisting}
// 二进制中 1 的个数
__builtin_popcount(x);     // unsigned int 版
__builtin_popcountll(x);   // unsigned long long 版

// 末尾 0 的个数（即最低位 1 的位置）
__builtin_ctz(x);          // x != 0 时才有意义
__builtin_ctzll(x);

// 前导 0 的个数
__builtin_clz(x);          // x != 0 时才有意义
__builtin_clzll(x);

// 奇偶性（1 的个数是否为奇数）
__builtin_parity(x);       // 奇数个 1 返回 1
__builtin_parityll(x);

// 常用推导
// 最高位 1 的位置 = 63 - __builtin_clzll(x)   (x > 0)
// 最低位 1 的位置 = __builtin_ctzll(x)         (x > 0)
// 向上取 2 的幂 = 1LL << (64 - __builtin_clzll(x - 1))  (x > 1)
// 向下取 2 的幂 = 1LL << (63 - __builtin_clzll(x))      (x > 0)
// log2(x) 下取整 = 63 - __builtin_clzll(x)    (x > 0)
\end{lstlisting}

\begin{warn}
\begin{itemize}
  \item \lstinline|__builtin_clz(0)| 和 \lstinline|__builtin_ctz(0)| 是未定义行为！使用前必须判断 $x \neq 0$。
  \item \lstinline|int| 类型用不带 \lstinline|ll| 的版本，\lstinline|long long| 类型必须用 \lstinline|ll| 版本，否则结果错误。
\end{itemize}
\end{warn}

\section{C++ 常用函数速查}

\subsection{数值函数}

\begin{lstlisting}
// <algorithm> 头文件
min(a, b);  max(a, b);        // 最小值/最大值
min({a, b, c, d});            // 多个值取最小（C++11，初始化列表）
max({a, b, c, d});
swap(a, b);                   // 交换两个值
clamp(x, lo, hi);             // 限制 x 在 [lo, hi] 范围内（C++17）

// <numeric> 头文件
accumulate(a.begin(), a.end(), 0LL);      // 累加（注意初始值类型！）
accumulate(a.begin(), a.end(), 1LL, multiplies<>()); // 累乘
iota(a.begin(), a.end(), 0);              // 填充 0, 1, 2, ...
partial_sum(a.begin(), a.end(), pre.begin()); // 前缀和
adjacent_difference(a.begin(), a.end(), d.begin()); // 差分
inner_product(a.begin(), a.end(), b.begin(), 0LL);  // 内积
gcd(a, b);  lcm(a, b);       // 最大公约数/最小公倍数（C++17）
midpoint(a, b);               // (a+b)/2 不溢出（C++20）
reduce(a.begin(), a.end());   // 并行累加（C++17，无序）

// <cmath> 头文件
abs(x);                       // 绝对值（整数用 abs，浮点也可以）
ceil(x);  floor(x);  round(x); // 上取整/下取整/四舍五入
sqrt(x);  cbrt(x);            // 平方根/立方根
log(x);   log2(x);  log10(x); // 自然对数/以2为底/以10为底
pow(x, y);                    // x 的 y 次方（浮点，整数慎用！）
fmod(x, y);                   // 浮点取模
hypot(x, y);                  // sqrt(x*x + y*y)，不溢出
\end{lstlisting}

\begin{warn}
\begin{itemize}
  \item \lstinline|pow(x, y)| 返回浮点数，\lstinline|pow(5, 2)| 可能返回 24.999...，赋给 int 变成 24！整数幂次用快速幂。
  \item 整数除法上取整：$\lceil a/b \rceil$ = \lstinline|(a + b - 1) / b|（$a, b > 0$ 时）。不要用 \lstinline|ceil((double)a / b)|，浮点不精确。
  \item \lstinline|abs(INT_MIN)| 溢出！\lstinline|INT_MIN| 的绝对值超过 \lstinline|INT_MAX|。
\end{itemize}
\end{warn}

\subsection{区间操作函数}

\begin{lstlisting}
// 排序
sort(a.begin(), a.end());                    // 升序
sort(a.begin(), a.end(), greater<int>());    // 降序
stable_sort(a.begin(), a.end());             // 稳定排序
partial_sort(a.begin(), a.begin()+k, a.end()); // 只排前 k 小
nth_element(a.begin(), a.begin()+k, a.end());  // 第 k 小放到位置 k，O(n)

// 二分查找（要求有序）
lower_bound(a.begin(), a.end(), x);  // >= x 的第一个位置
upper_bound(a.begin(), a.end(), x);  // > x 的第一个位置
binary_search(a.begin(), a.end(), x); // 是否存在 x（bool）
equal_range(a.begin(), a.end(), x);  // 返回 {lower_bound, upper_bound} 的 pair

// 查找
find(a.begin(), a.end(), x);         // 找第一个 x，返回迭代器
count(a.begin(), a.end(), x);        // 统计 x 的个数
// 接口：判定 lambda(v)：判断当前数组元素 v 是否大于 0。
// 参数：当前被 count_if 判断的数组元素 v。
count_if(a.begin(), a.end(), [](int v){ return v > 0; }); // 条件统计
any_of(a.begin(), a.end(), pred);    // 是否存在满足条件的
all_of(a.begin(), a.end(), pred);    // 是否全部满足条件
none_of(a.begin(), a.end(), pred);   // 是否全不满足条件

// 最值
*max_element(a.begin(), a.end());     // 最大值
*min_element(a.begin(), a.end());     // 最小值
auto [mi, ma] = minmax_element(a.begin(), a.end()); // 同时求最小和最大的迭代器

// 去重（要求有序）
auto it = unique(a.begin(), a.end()); // 返回去重后的尾迭代器
a.erase(it, a.end());                // 真正删除多余元素

// 反转、旋转、填充
reverse(a.begin(), a.end());          // 反转
rotate(a.begin(), a.begin()+k, a.end()); // 左旋 k 位
fill(a.begin(), a.end(), val);        // 填充
fill_n(a.begin(), k, val);           // 填充前 k 个
// 接口：生成 lambda()：调用 rng() 生成当前要写入数组的随机值。
generate(a.begin(), a.end(), [&](){ return rng(); }); // 用函数生成

// 拷贝、变换
copy(a.begin(), a.end(), b.begin());  // 拷贝到 b
// 接口：变换 lambda(x)：返回输入元素 x 的两倍并写入目标数组。
// 参数：当前被 transform 变换的输入元素 x。
transform(a.begin(), a.end(), b.begin(), [](int x){ return x*2; }); // 变换
\end{lstlisting}

\begin{tip}
\begin{itemize}
  \item \lstinline|nth_element| 是 $O(n)$ 的，比 \lstinline|sort| 后取第 k 个快很多。常用于求中位数。
  \item \lstinline|partial_sort| 只排前 k 小，$O(n \log k)$，当 $k \ll n$ 时比完整排序快。
  \item \lstinline|rotate| 实现数组左旋非常方便，$O(n)$。
\end{itemize}
\end{tip}

\subsection{排列与组合}

\begin{lstlisting}
// 全排列
// next_permutation: 生成下一个字典序排列，到最后返回 false
// prev_permutation: 生成上一个字典序排列，到最前返回 false
// 变量：a=用于演示排列算法的整数序列 a。
vector<int> a = {1, 2, 3};
do {
    // 处理当前排列 a
} while (next_permutation(a.begin(), a.end()));
// 注意：要枚举所有排列，a 必须先排序为最小排列！

// 枚举子集（用 bitmask）
// 变量：mask=用二进制位编码的当前状态 mask。
for (int mask = 0; mask < (1 << n); mask++) {
    // mask 的第 i 位为 1 表示选第 i 个元素
    // 变量：i=当前子集掩码中正在检查是否被选中的元素位 i。
    for (int i = 0; i < n; i++)
        if (mask >> i & 1) { /* 选了第 i 个 */ }
}

// 枚举大小为 k 的子集（用 bitmask）
// 见位运算章节

// is_permutation: 判断两个序列是否互为排列
is_permutation(a.begin(), a.end(), b.begin()); // bool
\end{lstlisting}

\begin{warn}
\begin{itemize}
  \item \lstinline|next_permutation| 会修改原数组！如果要枚举所有排列，初始必须是最小排列（升序）。
  \item 全排列总数是 $n!$，$n \ge 12$ 时 $n! > 10^8$，注意复杂度。
  \item \lstinline|prev_permutation| 生成上一个排列，初始必须是最大排列（降序）才能枚举所有排列。
\end{itemize}
\end{warn}

\subsection{全排列函数详解（next\_permutation / prev\_permutation）}

\begin{note}
C++ STL 提供的 \lstinline|next_permutation| 和 \lstinline|prev_permutation| 是生成排列的核心工具。它们\textbf{原地修改}序列，将其变换为字典序的下一个/上一个排列，返回 \lstinline|bool| 表示是否成功（到达最后/最前时返回 \lstinline|false| 并回绕为最小/最大排列）。时间复杂度单次 $O(n)$，枚举所有排列总计 $O(n! \cdot n)$。
\end{note}

\begin{lstlisting}
// 函数签名
bool next_permutation(BidirIt first, BidirIt last);
bool next_permutation(BidirIt first, BidirIt last, Compare comp);
bool prev_permutation(BidirIt first, BidirIt last);
bool prev_permutation(BidirIt first, BidirIt last, Compare comp);
// 返回值：成功生成下一个/上一个排列返回 true
//         已经是最大/最小排列时回绕并返回 false

// ===== 用法一：枚举所有排列 =====
// 必须先排序为最小排列（升序），然后用 do-while 循环
// 变量：a=用于演示排列算法的整数序列 a。
vector<int> a = {1, 2, 3};
sort(a.begin(), a.end()); // 必须先升序排列！
do {
    // 处理当前排列，例如输出
    // 变量：x=从 a 枚举的元素值 x。
    for (int x : a) cout << x << " ";
    cout << endl;
} while (next_permutation(a.begin(), a.end()));
// 输出：1 2 3 / 1 3 2 / 2 1 3 / 2 3 1 / 3 1 2 / 3 2 1

// ===== 用法二：有重复元素的排列（自动去重） =====
// next_permutation 天然支持重复元素，不会生成重复排列
// 变量：a=用于演示排列算法的整数序列 a。
vector<int> a = {1, 1, 2};
sort(a.begin(), a.end());
do {
    // 变量：x=从 a 枚举的元素值 x。
    for (int x : a) cout << x << " ";
    cout << endl;
} while (next_permutation(a.begin(), a.end()));
// 输出：1 1 2 / 1 2 1 / 2 1 1（只有 3 种，不是 3!=6 种）

// ===== 用法三：只排列部分元素（前 k 个的全排列） =====
// 枚举 n 个元素中选 k 个的全排列（P(n,k)）
// 变量：n=被选元素全集的大小 n；k=需要从 n 个元素中选取的元素个数 k。
int n = 4, k = 2;
// 变量：a=用于演示排列算法的整数序列 a。
vector<int> a = {1, 2, 3, 4};
sort(a.begin(), a.end());
do {
    // 只看前 k 个元素，它们是一种选法的排列
    // 变量：i=当前输出所选组合中第 i 个元素的下标 i。
    for (int i = 0; i < k; i++) cout << a[i] << " ";
    cout << endl;
    // 将后面的部分反转，使 next_permutation 跳到下一组前 k 个
    reverse(a.begin() + k, a.end());
} while (next_permutation(a.begin(), a.end()));

// ===== 用法四：求当前排列之后的第 k 个排列 =====
// 变量：a=用于演示排列算法的整数序列 a。
vector<int> a = {1, 2, 3, 4};
// 变量：k=限制输出的前 k 个排列数量 k。
int k = 5;
while (k-- && next_permutation(a.begin(), a.end()));
// a 变成第 5 个后继排列

// ===== 用法五：对字符串使用 =====
// 变量：s=从标准输入读取、供当前字符串算法处理的字符串 s。
string s = "abc";
sort(s.begin(), s.end());
do {
    cout << s << endl;
} while (next_permutation(s.begin(), s.end()));
// 输出：abc / acb / bac / bca / cab / cba

// ===== 用法六：自定义比较函数（降序排列枚举） =====
vector<int> a = {3, 2, 1}; // 按 greater 的最小排列
do {
    // 变量：x=从 a 枚举的元素值 x。
    for (int x : a) cout << x << " ";
    cout << endl;
} while (next_permutation(a.begin(), a.end(), greater<int>()));
// 枚举所有"降序字典序"的排列
\end{lstlisting}

\begin{warn}
\begin{itemize}
  \item \textbf{初始状态很重要}：用 \lstinline|next_permutation| 枚举所有排列，初始必须是升序（字典序最小）；用 \lstinline|prev_permutation| 则初始必须是降序（字典序最大）。否则只能枚举到从当前位置到末尾的部分排列。
  \item \textbf{回绕行为}：当 \lstinline|next_permutation| 在最大排列上调用时，会将数组变回最小排列并返回 \lstinline|false|。这就是 \lstinline|do-while| 的终止条件——回绕后 \lstinline|false| 退出循环。
  \item \textbf{复杂度}：$n!$ 增长极快，$n = 8$ 时 $n! = 40320$，$n = 10$ 时 $n! = 3628800$，$n = 12$ 时 $n! \approx 4.8 \times 10^8$。竞赛中 $n \le 8$ 比较安全，$n \le 10$ 勉强可行。
  \item \textbf{不要用 while}：写 \lstinline|while (next_permutation(...))| 会跳过初始排列！必须用 \lstinline|do { ... } while (next_permutation(...))|。
\end{itemize}
\end{warn}

\begin{tip}
\begin{itemize}
  \item \textbf{判断是否为排列}：用 \lstinline|is_permutation(a.begin(), a.end(), b.begin())| 判断 $a$ 和 $b$ 是否互为排列。
  \item \textbf{next\_permutation 的原理}：从右往左找第一个 $a[i] < a[i+1]$ 的位置，再从右往左找第一个 $a[j] > a[i]$ 的位置，交换 $a[i]$ 和 $a[j]$，最后反转 $a[i+1 \ldots]$。
  \item \textbf{求排列数/排名}：需要知道某排列是第几个，用康托展开（见下节）。
  \item \textbf{竞赛经典场景}：暴力枚举所有排列检查是否满足条件（如 N 皇后的暴力解、TSP 暴力等）。
\end{itemize}
\end{tip}

\subsubsection{DFS 递归生成全排列}

\begin{note}
当需要在生成过程中剪枝（如 N 皇后），或需要更灵活的控制（如只要部分元素的排列），用 DFS 递归比 \lstinline|next_permutation| 更方便。
\end{note}

\begin{lstlisting}
// 生成 {1, 2, ..., n} 的全排列（不重复元素）
// 变量：path=排列 DFS 当前已经选定的前缀 path。
vector<int> path;
// 变量：used=记录每个对象是否已被使用的标记 used。
vector<bool> used;

// 接口：dfs(n)：枚举并逐个输出 1..n 的全部排列；不返回值。
// 参数：待生成排列的长度 n。
void dfs(int n) {
    if ((int)path.size() == n) {
        // 得到一个排列 path，处理它
        // 变量：x=从 path 枚举的元素值 x。
        for (int x : path) cout << x << " ";
        cout << endl;
        return;
    }
    // 变量：i=无重复全排列 DFS 当前尝试放入路径的候选值 i。
    for (int i = 1; i <= n; i++) {
        if (used[i]) continue;
        path.push_back(i);
        used[i] = true;
        dfs(n);
        path.pop_back();
        used[i] = false;
    }
}
// 调用：used.assign(n + 1, false); dfs(n);

// 有重复元素的 DFS 全排列（去重版）
// nums 必须先排序！
vector<int> nums; // 已排序
// 变量：path2=含重复元素时的当前 DFS 排列 path2。
vector<int> path2;
// 变量：used2=含重复元素时每个下标是否已被当前排列使用。
vector<bool> used2;

// 接口：dfs2(n)：枚举已排序 nums 的所有不重复全排列并逐个输出；不返回值。
// 参数：待生成排列的长度 n。
void dfs2(int n) {
    if ((int)path2.size() == n) {
        // 变量：x=从 path2 枚举的元素值 x。
        for (int x : path2) cout << x << " ";
        cout << endl;
        return;
    }
    // 变量：i=含重复元素全排列 DFS 当前尝试使用的 nums 下标 i。
    for (int i = 0; i < n; i++) {
        if (used2[i]) continue;
        // 去重：同一层中，相同值只选第一个未用的
        if (i > 0 && nums[i] == nums[i-1] && !used2[i-1]) continue;
        path2.push_back(nums[i]);
        used2[i] = true;
        dfs2(n);
        path2.pop_back();
        used2[i] = false;
    }
}
// 调用：sort(nums.begin(), nums.end()); used2.assign(n, false); dfs2(n);
\end{lstlisting}

\begin{tip}
\begin{itemize}
  \item DFS 版本可以方便地在递归中\textbf{剪枝}，例如 N 皇后问题在放置每个皇后时检查冲突，大幅减少搜索量。
  \item DFS 生成\textbf{部分排列}（选 $k$ 个）：只需把终止条件改为 \lstinline|path.size() == k|。
  \item 去重的关键条件 \lstinline|!used2[i-1]| 表示：如果前一个相同元素没有被选（说明是同层回溯回来的），则跳过当前元素，避免重复。
\end{itemize}
\end{tip}

\subsection{排列序（康托展开与逆康托展开）}

\begin{note}
康托展开：将一个排列映射为它在所有排列中的字典序排名（0-indexed）。逆康托展开：由排名恢复排列。时间复杂度均为 $O(n^2)$（用树状数组可优化到 $O(n \log n)$）。
\end{note}

\begin{lstlisting}
// 康托展开：求排列 p 在所有 n! 个排列中的排名（0-indexed）
// p 是 {0, 1, ..., n-1} 的排列
// 接口：cantor(p)：返回排列 p 的 0-based Cantor 排名；返回类型 int。
// 参数：待计算排名的 0-based 排列 p。
int cantor(vector<int>& p) {
    // 变量：n=容器 p 的元素数量 n。
    int n = p.size();
    // 变量：rank=当前排列的 0-based Cantor 排名 rank。
    int rank = 0;
    // 变量：fact=阶乘表 fact；fact[i]=i!。
    vector<int> fact(n, 1);
    // 变量：i=Cantor 排名使用的阶乘表下标 i。
    for (int i = 1; i < n; i++) fact[i] = fact[i-1] * i;
    // 变量：i=Cantor 排名中当前计算较小后继数量的排列位置 i。
    for (int i = 0; i < n; i++) {
        int cnt = 0; // 统计 p[i] 后面比 p[i] 小的个数
        // 变量：j=Cantor 排名中位于 i 之后、用于统计较小元素的位置 j。
        for (int j = i + 1; j < n; j++)
            if (p[j] < p[i]) cnt++;
        rank += cnt * fact[n - 1 - i];
    }
    return rank;
}

// 逆康托展开：由排名 rank 恢复排列（0-indexed）
// 返回 {0, 1, ..., n-1} 的第 rank 个排列
// 接口：inv_cantor(rank, n)：由 0-based 排名还原 1..n 的排列；返回类型 vector<int>。
// 参数：待还原的 0-based Cantor 排名 rank；待还原排列的长度 n。
vector<int> inv_cantor(int rank, int n) {
    // 变量：fact=阶乘表 fact；fact[i]=i!。
    vector<int> fact(n, 1);
    // 变量：i=Cantor 排名使用的阶乘表下标 i。
    for (int i = 1; i < n; i++) fact[i] = fact[i-1] * i;
    vector<int> avail; // 可用元素
    // 变量：i=逆 Cantor 展开初始化可用值集合的元素值 i。
    for (int i = 0; i < n; i++) avail.push_back(i);
    // 变量：res=逆 Cantor 展开逐位构造的排列 res。
    vector<int> res;
    // 变量：i=逆 Cantor 展开中当前确定输出值的排列位置 i。
    for (int i = 0; i < n; i++) {
        // 变量：idx=当前访问容器的 0-based 下标 idx。
        int idx = rank / fact[n - 1 - i];
        rank %= fact[n - 1 - i];
        res.push_back(avail[idx]);
        avail.erase(avail.begin() + idx);
    }
    return res;
}
\end{lstlisting}

\begin{tip}
\begin{itemize}
  \item 如果排列是 $\{1, 2, \ldots, n\}$（1-indexed），只需将输入减 1 再调用，或将 avail 改为 $1 \sim n$。
  \item 用树状数组替代内层循环可以优化到 $O(n \log n)$（统计前缀中比当前小的个数）。
  \item 康托展开在排列哈希、状态压缩（如 8-puzzle）中非常有用。
  \item $n \le 20$ 时 $n!$ 超过 \lstinline|long long| 范围，需要用高精度或只处理小规模。$n \le 12$ 时 \lstinline|int| 足够，$n \le 20$ 用 \lstinline|long long|。
\end{itemize}
\end{tip}

\subsection{字符串处理函数}

\begin{lstlisting}
// <string> 头文件
s.substr(pos, len);            // 子串
s.find("abc");                 // 查找子串位置，找不到返回 string::npos
s.rfind("abc");                // 从右往左找
s.find_first_of("aeiou");     // 找第一个属于给定字符集的字符
s.find_last_of("aeiou");      // 找最后一个属于给定字符集的字符
s.find_first_not_of("01");    // 找第一个不在给定字符集中的字符
s.replace(pos, len, "new");   // 替换子串
s.insert(pos, "abc");         // 在 pos 处插入
s.erase(pos, len);            // 删除从 pos 开始的 len 个字符
s.starts_with("pre");         // 是否以 "pre" 开头（C++20）
s.ends_with("suf");           // 是否以 "suf" 结尾（C++20）
s.contains("sub");            // 是否包含子串（C++23）

// 类型转换
to_string(123);                // int -> string
to_string(3.14);               // double -> string
stoi("123");                   // string -> int
stoll("123456789012");         // string -> long long
stod("3.14");                  // string -> double

// <cctype> 字符分类（传入 char）
isalpha(c);   // 是否是字母
isdigit(c);   // 是否是数字
isalnum(c);   // 是否是字母或数字
isupper(c);   // 是否是大写字母
islower(c);   // 是否是小写字母
isspace(c);   // 是否是空白字符
toupper(c);   // 转大写
tolower(c);   // 转小写

// 整行读入
getline(cin, s);  // 读一整行（注意前面 cin>>n 后要 cin.ignore()）

// 字符串流（<sstream>）
stringstream ss("1 2 3");
// 变量：a=字符串解析得到的第一个整数 a；b=字符串解析得到的第二个整数 b；c=字符串解析得到的第三个整数 c。
int a, b, c;
ss >> a >> b >> c;  // 按空格分割解析
\end{lstlisting}

\begin{warn}
\begin{itemize}
  \item \lstinline|s.find()| 找不到时返回 \lstinline|string::npos|（一个极大的无符号数），不是 $-1$！判断时用 \lstinline|if (s.find(x) != string::npos)|。
  \item \lstinline|stoi| 对超出 int 范围的字符串会抛异常，大数用 \lstinline|stoll|。
  \item \lstinline|getline| 前如果用了 \lstinline|cin >>|，必须先 \lstinline|cin.ignore()| 消耗掉残留的换行符，否则读到空行。
\end{itemize}
\end{warn}

\subsection{其他实用工具}

\begin{lstlisting}
// 二分答案的 lambda 写法
// 接口：check(mid)：示范二分答案判定接口；调用者必须把函数体替换为题目的 mid 可行性条件。
// 参数：当前待判定的二分答案 mid。
auto check = [&](int mid) -> bool {
    // 判断 mid 是否可行
    return true;
};

// 自定义哈希（pair 做 key）
struct pair_hash {
    // 接口：operator()：返回键值的 size_t 哈希值，供 unordered_map/unordered_set 去重；返回类型 size_t。
    size_t operator()(const pair<int,int>& p) const {
        return hash<long long>()(((long long)p.first << 32) | p.second);
    }
};
// 变量：mp=使用 pair_hash 把整数对键映射到 int 值的哈希表 mp。
unordered_map<pair<int,int>, int, pair_hash> mp;
// 变量：st=用于按 pair 键去重的无序集合 st。
unordered_set<pair<int,int>, pair_hash> st;

// 离散化
vector<int> vals(a.begin(), a.end()); // 复制原数组
sort(vals.begin(), vals.end());
vals.erase(unique(vals.begin(), vals.end()), vals.end());
// 映射：原值 -> 离散化后的编号
// 变量：i=坐标离散化时当前回写压缩编号的数组下标 i。
for (int i = 0; i < n; i++)
    a[i] = lower_bound(vals.begin(), vals.end(), a[i]) - vals.begin();
// 反映射：vals[a[i]] 即为原值

// iota 快速填充连续序列
// 变量：idx=当前访问容器的 0-based 下标 idx。
vector<int> idx(n);
iota(idx.begin(), idx.end(), 0); // idx = {0, 1, 2, ..., n-1}
// 按其他数组的值排序下标（间接排序）
// 接口：比较器 lambda(i, j)：按 a[i]、a[j] 的值升序排列下标。
// 参数：第一个候选下标 i；第二个候选下标 j。
sort(idx.begin(), idx.end(), [&](int i, int j) {
    return a[i] < a[j];
});

// merge 合并两个有序序列
merge(a.begin(), a.end(), b.begin(), b.end(), c.begin());

// set_intersection / set_union / set_difference（要求有序）
// 变量：res=当前筛选操作收集的整数结果 res。
vector<int> res;
set_intersection(a.begin(), a.end(), b.begin(), b.end(), back_inserter(res));
// 类似地：set_union, set_difference, set_symmetric_difference
\end{lstlisting}

\begin{tip}
\begin{itemize}
  \item 离散化是竞赛高频操作，用于将大值域压缩到 $[0, n)$，配合树状数组/线段树使用。
  \item 间接排序（按下标排序）可以保留原数组不变，同时知道排序后每个元素的原始位置。
  \item \lstinline|merge| 和集合运算函数要求输入有序，输出也有序。
\end{itemize}
\end{tip}

% ==========================================================
\chapter{位运算}
% ==========================================================

\section{基本性质}

\begin{note}
\begin{itemize}
  \item 与 \lstinline|&|：同 1 为 1。交换律、结合律。$a \& 0 = 0$，$a \& a = a$，$a \& (\sim a) = 0$。
  \item 或 \lstinline{|}：有 1 为 1。交换律、结合律。$a | 0 = a$，$a | a = a$，$a | (\sim a) = -1$（全 1）。
  \item 异或 \lstinline|^|：不同为 1。交换律、结合律。$a \oplus 0 = a$，$a \oplus a = 0$，$a \oplus (\sim a) = -1$。
  \item 取反 \lstinline|~|：$\sim a = -a - 1$（补码性质）。
\end{itemize}
\end{note}

\subsection{分配律与吸收律}

\begin{note}
\begin{itemize}
  \item 与对或的分配律：$a \& (b | c) = (a \& b) | (a \& c)$
  \item 或对与的分配律：$a | (b \& c) = (a | b) \& (a | c)$
  \item 异或对与的分配律：$a \& (b \oplus c) = (a \& b) \oplus (a \& c)$
  \item 吸收律：$a \& (a | b) = a$，$a | (a \& b) = a$
  \item 德摩根律：$\sim(a \& b) = (\sim a) | (\sim b)$，$\sim(a | b) = (\sim a) \& (\sim b)$
\end{itemize}
\end{note}

\subsection{异或的重要性质}

\begin{note}
\begin{itemize}
  \item 自逆性：$a \oplus b \oplus b = a$（异或同一个数两次等于没异或，用于加密/交换）。
  \item 无进位加法：异或等价于不考虑进位的二进制加法。
  \item 前缀异或：$a_l \oplus a_{l+1} \oplus \cdots \oplus a_r = \text{pre}[r] \oplus \text{pre}[l-1]$，类比前缀和。
  \item 交换两数（不用临时变量）：\lstinline|a ^= b; b ^= a; a ^= b;|（但实际竞赛中直接用 swap 更安全）。
  \item 线性基：$n$ 个数的所有异或组合可以用至多 $\log V$ 个基底表示。
\end{itemize}
\end{note}

\subsection{$O(1)$ 求连续自然数的异或和}

\begin{note}
求 $1 \oplus 2 \oplus \cdots \oplus m$ 是一个经典结论。令 $S(m) = \bigoplus_{i=1}^{m} i$，则结果只取决于 $m \bmod 4$：
\begin{itemize}
  \item $m \bmod 4 = 0$ 时，$S(m) = m$
  \item $m \bmod 4 = 1$ 时，$S(m) = 1$
  \item $m \bmod 4 = 2$ 时，$S(m) = m + 1$
  \item $m \bmod 4 = 3$ 时，$S(m) = 0$
\end{itemize}

\textbf{为什么有这个结论}：核心引理是\textbf{每 4 个对齐到 4 的倍数的连续整数，异或起来等于 0}。即对任意 $k \equiv 0 \pmod 4$：
$$k \oplus (k+1) \oplus (k+2) \oplus (k+3) = 0.$$
原因：这 4 个数的\textbf{低 2 位}恰好取遍 $00,01,10,11$，异或为 $0$；而\textbf{高位}（第 2 位及以上）四个数完全相同，相同的数异或偶数次也为 $0$。所以整块异或为 $0$。

把 $S(m)$ 写成 $0 \oplus 1 \oplus \cdots \oplus m$（多异或一个 $0$ 不影响结果），再从 $0$ 开始每 4 个分一组：$(0\,1\,2\,3)(4\,5\,6\,7)\cdots$，每个\textbf{完整块}都异或为 $0$，于是只剩下结尾不满一块的零头：
\begin{itemize}
  \item $m \equiv 3$：恰好若干完整块，全抵消，$S(m) = 0$。
  \item $m \equiv 0$：到 $m-1$（$\equiv 3$）为完整块抵消，再异或上多出来的 $m$，$S(m) = m$。
  \item $m \equiv 1$：$S(m) = S(m-1) \oplus m = (m-1) \oplus m$。此时 $m$ 为奇数，$m-1$ 与 $m$ 只差最低位，故 $(m-1)\oplus m = 1$。
  \item $m \equiv 2$：$S(m) = S(m-1) \oplus m = 1 \oplus m$。此时 $m$ 为偶数（最低位为 $0$），故 $1 \oplus m = m + 1$。
\end{itemize}

\textbf{推论（区间异或和）}：利用前缀异或，任意区间 $\bigoplus_{i=l}^{r} i = S(r) \oplus S(l-1)$，同样 $O(1)$。
\end{note}

\begin{lstlisting}
// 1 ^ 2 ^ ... ^ m，O(1)
// 接口：xor_n(m)：返回整数闭区间 [1,m] 的异或和；返回类型 int。
// 参数：异或前缀的右端点 m。
int xor_n(int m) {
    switch (m & 3) {        // m & 3 等价于 m % 4
        case 0: return m;
        case 1: return 1;
        case 2: return m + 1;
        default: return 0;  // case 3
    }
}

// 区间异或和 l ^ (l+1) ^ ... ^ r，O(1)
// 接口：xor_range(l, r)：返回闭区间 [l,r] 内所有整数的异或和；返回类型 int。
// 参数：闭区间左端点 l；闭区间右端点 r。
int xor_range(int l, int r) {
    return xor_n(r) ^ xor_n(l - 1);
}
\end{lstlisting}

\subsection{与/或的单调性}

\begin{note}
\begin{itemize}
  \item 与运算越与越小：$a \& b \le \min(a, b)$。连续与一串数，结果单调不增。
  \item 或运算越或越大：$a | b \ge \max(a, b)$。连续或一串数，结果单调不减。
  \item 推论：子数组的 AND 值最多有 $O(\log V)$ 种（每次变化至少少一个 1）；OR 同理。
  \item 应用：枚举子数组的 AND/OR 值时，可以利用单调性 + 双指针/二分优化。
\end{itemize}
\end{note}

\subsection{常用恒等式速查}

\begin{note}
\begin{itemize}
  \item $a \& (a-1)$：去掉 $a$ 最低位的 1。应用：统计 1 的个数、判断 2 的幂。
  \item $a \& (-a)$：取出 $a$ 最低位的 1（lowbit）。应用：树状数组。
  \item $a | (a+1)$：将 $a$ 最低位的 0 置为 1。
  \item $a \& (a+1)$：将 $a$ 末尾连续的 1 全部清零。
  \item $a | (a-1)$：将 $a$ 末尾连续的 0 全部置 1。
  \item $(a \oplus b) | (a \& b) = a | b$
  \item $(a \oplus b) \& (a \& b) = 0$（异或和与互斥）
  \item $a + b = (a \oplus b) + 2 \cdot (a \& b)$（加法 = 无进位和 + 进位）
  \item $a + b = (a | b) + (a \& b)$（另一个加法分解）
\end{itemize}
\end{note}

\section{竞赛常用性质科普}

\subsection{按位独立性（Bitwise Independence）}

\begin{note}
位运算中，AND / OR / XOR 对每一位的运算是独立的——第 $i$ 位的结果只取决于所有操作数的第 $i$ 位。

\textbf{竞赛意义}：遇到"对一组数做位运算后求最优值"的题目，可以逐位贪心或逐位 DP，把一个 $\log V$ 位的问题拆成 $\log V$ 个独立的 0/1 问题。

典型应用：
\begin{itemize}
  \item 求"选若干数使 OR 最大"——从高位到低位贪心，每一位独立决策。
  \item 求"所有子数组 AND 之和"——对每一位分别统计有多少子数组该位为 1。
  \item 位运算卷积（AND/OR/XOR 卷积）的正确性也依赖于按位独立性。
\end{itemize}
\end{note}

\subsection{异或与线性代数}

\begin{note}
异或运算构成 $\text{GF}(2)$（二元域）上的加法，AND 构成乘法。因此：
\begin{itemize}
  \item $n$ 个数的所有异或组合构成一个线性空间，维度 $\le \log V$。
  \item 线性基（Gaussian Elimination on GF(2)）可以在 $O(n \log V)$ 时间内求出基底。
  \item 线性基支持：求最大/最小异或值、判断某个值能否被异或出、求第 $k$ 小异或值。
  \item 异或方程组可以用高斯消元求解（开关灯问题、翻转问题）。
\end{itemize}
\end{note}

\subsection{位运算与集合论的对应}

\begin{note}
将整数看作集合（第 $i$ 位为 1 表示元素 $i$ 在集合中），位运算对应集合运算：
\begin{itemize}
  \item $a \& b$ = $A \cap B$（交集）
  \item $a | b$ = $A \cup B$（并集）
  \item $a \oplus b$ = $A \triangle B$（对称差）
  \item $\sim a$ = $\overline{A}$（补集）
  \item $a \& (\sim b)$ = $A \setminus B$（差集）
  \item $(a \& b) == a$ 等价于 $A \subseteq B$（子集判断）
\end{itemize}

\textbf{竞赛意义}：状压 DP 中，集合操作全部用位运算实现，理解这层对应关系可以快速写出状态转移。
\end{note}

\subsection{popcount 的性质与应用}

\begin{note}
$\text{popcount}(x)$ 表示 $x$ 二进制中 1 的个数。常用性质：
\begin{itemize}
  \item $\text{popcount}(a \oplus b)$ = $a$ 和 $b$ 不同位的个数（汉明距离）。
  \item $\text{popcount}(a \& b)$ = $a$ 和 $b$ 同时为 1 的位数。
  \item $\text{popcount}(a | b) = \text{popcount}(a) + \text{popcount}(b) - \text{popcount}(a \& b)$（容斥）。
  \item $a + b = 2 \cdot (a \& b) + (a \oplus b)$，因此 $\text{popcount}(a) + \text{popcount}(b) = \text{popcount}(a \oplus b) + 2 \cdot \text{popcount}(a \& b)$。
\end{itemize}

\textbf{竞赛应用}：
\begin{itemize}
  \item 状压 DP 中按 popcount 分层枚举（先枚举选 1 个的状态，再选 2 个……）。
  \item 判断两个状态的"距离"用汉明距离。
  \item 用 \lstinline|__builtin_popcountll(x)| 获取，$O(1)$。
\end{itemize}
\end{note}

\subsection{高位/低位操作}

\begin{note}
\begin{itemize}
  \item 最高位位置：$\lfloor \log_2 x \rfloor$ = \lstinline|63 - __builtin_clzll(x)|（$x > 0$）。
  \item 最低位位置：\lstinline|__builtin_ctzll(x)|（$x > 0$）。
  \item 只保留最高位的 1：\lstinline|1LL << (63 - __builtin_clzll(x))|。
  \item 只保留最低位的 1：\lstinline|x & (-x)|（lowbit）。
  \item 将最高位以下全部填 1：先求最高位位置 $k$，然后 \lstinline|(1LL << (k+1)) - 1|。
\end{itemize}

\textbf{竞赛意义}：
\begin{itemize}
  \item 贪心从高位到低位确定答案时，需要快速定位最高位。
  \item 树状数组的核心操作就是 lowbit。
  \item 求"不超过 $x$ 的最大 2 的幂"等问题直接用最高位操作。
\end{itemize}
\end{note}

\subsection{位运算与不等式}

\begin{note}
\begin{itemize}
  \item $a \& b \le \min(a, b)$，$a | b \ge \max(a, b)$。
  \item $a \oplus b \le a + b$（异或不会产生进位，加法会）。
  \item $a + b = (a | b) + (a \& b) \ge a | b$。
  \item $(a \oplus b) + (a \& b) \cdot 2 = a + b$。
  \item 若 $a \& b = 0$（无公共位），则 $a + b = a | b = a \oplus b$。
\end{itemize}

\textbf{竞赛意义}：
\begin{itemize}
  \item "选两个数使 AND 最大"——从高位贪心，尽量让高位同时为 1。
  \item "选两个数使 XOR 最大"——用异或字典树从高位贪心取不同位。
  \item $a \& b = 0$ 时三种运算等价，常用于状压 DP 中"两个集合不相交"的条件。
\end{itemize}
\end{note}

\subsection{SOS DP（子集和 / 超集和）}

\begin{note}
给定数组 $f[mask]$，求所有子集的和：$g[mask] = \sum_{sub \subseteq mask} f[sub]$。

暴力枚举子集是 $O(3^n)$，SOS DP 可以做到 $O(n \cdot 2^n)$：
\begin{lstlisting}
// 数值域：0<=n<31，且每个子集和都必须可由 f 的元素类型表示；否则把 f 改为 __int128。
// 变量：i=SOS DP 当前向所有掩码传播贡献的二进制位 i。
for (int i = 0; i < n; i++)
    // 变量：mask=用二进制位编码的当前状态 mask。
    for (int mask = 0; mask < (1 << n); mask++)
        if (mask >> i & 1)
            f[mask] += f[mask ^ (1 << i)];
\end{lstlisting}

类似地，超集和（$g[mask] = \sum_{mask \subseteq sup} f[sup]$）只需把条件改为 \lstinline|if (!(mask >> i & 1))|。

\textbf{竞赛应用}：
\begin{itemize}
  \item 求"有多少对 $(i,j)$ 满足 $a_i \& a_j = a_i$（即 $a_i$ 是 $a_j$ 的子集）"。
  \item 位运算卷积（AND/OR 卷积）的核心变换。
  \item 容斥原理在位运算上的高效实现。
\end{itemize}
\end{note}

\section{常用技巧}

\begin{lstlisting}
// lowbit：取出最低位的 1
// 接口：lowbit(x)：返回 x 的二进制最低位 1 所代表的权值 x&-x；返回类型 int。
// 参数：要求最低位 1 权值的整数 x。
int lowbit(int x) { return x & (-x); }

// 判断是否是 2 的幂
// 接口：is_pow2(x)：判断正整数 x 是否恰好是 2 的幂；x>0 且二进制中仅有一个 1 时返回 true。
// 参数：待判定是否为 2 的幂的正整数 x。
bool is_pow2(int x) { return x > 0 && (x & (x - 1)) == 0; }

// 枚举子集（枚举 mask 的所有非空子集）
// 变量：sub=当前枚举的子掩码 sub。
for (int sub = mask; sub; sub = (sub - 1) & mask) {
    // sub 是 mask 的一个子集
}

// 枚举大小为 k 的子集（从 n 个元素中选 k 个）
// 变量：s=从前 k 个二进制位全为 1 开始枚举的子集掩码 s。
int s = (1 << k) - 1;
while (s < (1 << n)) {
    // 处理 s
    // 变量：lb=Gosper 枚举中当前掩码 s 的最低位 1 对应权值 lb。
    int lb = s & (-s);
    // 变量：r=Gosper 枚举中大于 s 的下一个同 popcount 掩码构造量 r。
    int r = s + lb;
    s = (((r ^ s) >> 2) / lb) | r;
}
\end{lstlisting}

\begin{warn}
\begin{itemize}
  \item 位运算优先级极低！\lstinline|if (x & 1 == 1)| 实际是 \lstinline|x & (1 == 1)|，必须加括号 \lstinline|if ((x & 1) == 1)|。
  \item 左移 \lstinline|1 << i| 当 $i \ge 31$ 时溢出，必须写 \lstinline|1LL << i|。
\end{itemize}
\end{warn}

\inputnewboard{remake/large/01_位运算.tex}
\inputnewboard{remake/large/01_线性基_详解.tex}

% ==========================================================

\chapter{基础算法}
\inputnewboard{remake/large/02_基础.tex}

\chapter{数学}
\inputnewboard{remake/large/03_数学.tex}
\inputnewboard{remake/large/03_高精度.tex}
\inputnewboard{remake/large/03_数学_详解.tex}

\inputnewboard{remake/large/04_数据结构.tex}
\inputnewboard{remake/large/04_数据结构_详解.tex}

\inputnewboard{remake/large/05_字符串.tex}
\inputnewboard{remake/large/05_字符串_详解.tex}

\inputnewboard{remake/large/06_图论.tex}

\chapter{树上算法与高级结构}
\inputnewboard{remake/large/07_树上进阶.tex}

\chapter{动态规划}
\inputnewboard{remake/large/08_动态规划.tex}
\inputnewboard{remake/large/08_动态规划_详解.tex}

\inputnewboard{remake/large/09_博弈.tex}
\inputnewboard{remake/large/09_博弈_详解.tex}

\inputnewboard{remake/large/10_几何.tex}
\inputnewboard{remake/large/10_几何_详解.tex}
\inputnewboard{remake/large/11_针对性测试.tex}
\inputnewboard{remake/large/13_例题与测试.tex}

\chapter{对拍}
% ==========================================================

\begin{note}
对拍（Stress Testing）是竞赛中查错的核心手段：用随机数据生成器产生大量测试数据，分别用暴力程序和待测程序运行，比较输出是否一致。一旦发现不一致，即可定位 bug。

对拍三件套：
\begin{itemize}
  \item \textbf{gen.cpp}：随机数据生成器，输出一组合法的输入数据。
  \item \textbf{brute.cpp}：暴力解法（正确但慢），作为对照标准。
  \item \textbf{sol.cpp}：你的正式解法（快但可能有 bug）。
\end{itemize}
\end{note}

\section{对拍脚本}

\subsection{Linux / macOS（Bash）}

\begin{lstlisting}
#!/bin/bash
# 编译三个程序
g++ -O2 -o gen gen.cpp
g++ -O2 -o brute brute.cpp
g++ -O2 -o sol sol.cpp

for ((i = 1; ; i++)); do
    # 生成随机数据（用循环变量 i 做种子）
    ./gen $i > in.txt
    # 分别运行暴力和正解
    ./brute < in.txt > out_brute.txt
    ./sol   < in.txt > out_sol.txt
    # 比较输出
    if diff out_brute.txt out_sol.txt > /dev/null 2>&1; then
        echo "Test $i: AC"
    else
        echo "Test $i: WA !!!"
        echo "--- Input ---"
        cat in.txt
        echo "--- Brute ---"
        cat out_brute.txt
        echo "--- Sol ---"
        cat out_sol.txt
        break
    fi
done
\end{lstlisting}

\subsection{Windows（BAT）}

\begin{lstlisting}
@echo off
g++ -O2 -o gen.exe gen.cpp
g++ -O2 -o brute.exe brute.cpp
g++ -O2 -o sol.exe sol.cpp

set i=0
:loop
set /a i+=1
gen.exe %i% > in.txt
brute.exe < in.txt > out_brute.txt
sol.exe   < in.txt > out_sol.txt
fc out_brute.txt out_sol.txt > nul
if errorlevel 1 (
    echo Test %i%: WA !!!
    echo --- Input ---
    type in.txt
    echo --- Brute ---
    type out_brute.txt
    echo --- Sol ---
    type out_sol.txt
    pause
    goto :eof
)
echo Test %i%: AC
goto loop
\end{lstlisting}

\begin{warn}
\begin{itemize}
  \item 编译时加 \lstinline|-O2| 让暴力跑得更快，能测更多组。
  \item 如果正解有可能输出多种正确答案（如最优方案不唯一），不能直接 diff，需要写 checker 来判断正确性。
  \item 如果对拍跑很久都没问题，考虑增大数据范围或换生成策略（边界、特殊图等）。
\end{itemize}
\end{warn}

\section{随机数据生成器（gen.cpp）}

\begin{note}
生成器的核心原则：
\begin{itemize}
  \item 用命令行参数作为随机种子，保证每次运行的数据不同，且可复现。
  \item 生成的数据必须满足题目的输入格式和约束条件。
  \item 对拍时先用小数据（$n \le 10$），确认暴力能跑出来。
\end{itemize}
\end{note}

\begin{lstlisting}
#include <bits/stdc++.h>
using namespace std;

mt19937 rng; // 梅森旋转随机数引擎

// 生成 [lo, hi] 范围内的随机整数
// 接口：randint(lo, hi)：返回闭区间 [l,r] 内均匀随机整数；返回类型 int。
// 参数：均匀随机整数闭区间下界 lo；均匀随机整数闭区间上界 hi。
int randint(int lo, int hi) {
    return uniform_int_distribution<int>(lo, hi)(rng);
}

// 生成 [lo, hi] 范围内的随机 long long
// 接口：randll(lo, hi)：返回闭区间 [l,r] 内均匀随机 long long；返回类型 long long。
// 参数：均匀随机 long long 闭区间下界 lo；均匀随机 long long 闭区间上界 hi。
long long randll(long long lo, long long hi) {
    return uniform_int_distribution<long long>(lo, hi)(rng);
}

// 生成长度为 len 的随机小写字母字符串
// 接口：randstr(len)：生成并返回长度为 len 的均匀随机小写字母字符串；返回类型 string。
// 参数：目标字符串长度 len。
string randstr(int len) {
    // 变量：s=正在逐字符构造并由 randstr 返回的随机小写字符串 s。
    string s(len, ' ');
    for (char& c : s) c = 'a' + randint(0, 25);
    return s;
}

// 生成 n 个节点的随机树（输出 n-1 条边）
// 方法：每个新节点随机连接到已有节点
// 接口：rand_tree(n)：生成 n 个点的随机树边集；返回类型 vector<pair<int,int>>。
// 参数：随机树的顶点数 n。
vector<pair<int,int>> rand_tree(int n) {
    // 变量：edges=当前随机树已经生成的无向边集合 edges。
    vector<pair<int,int>> edges;
    // 变量：i=随机树生成中当前新加入的顶点编号 i。
    for (int i = 2; i <= n; i++) {
        // 变量：fa=当前节点的父节点 fa。
        int fa = randint(1, i - 1);
        edges.push_back({fa, i});
    }
    // 打乱边的顺序和方向
    shuffle(edges.begin(), edges.end(), rng);
    for (auto& [u, v] : edges)
        if (randint(0, 1)) swap(u, v);
    return edges;
}

// 生成 n 个节点 m 条边的随机简单图（无重边无自环）
// 接口：rand_graph(n, m)：生成 n 点 m 边随机简单图；返回类型 vector<pair<int,int>>。
// 参数：随机图的顶点数 n；随机图的边数 m。
vector<pair<int,int>> rand_graph(int n, int m) {
    // 变量：used=记录每个对象是否已被使用的标记 used。
    set<pair<int,int>> used;
    // 变量：edges=当前随机简单图已经生成的无向边集合 edges。
    vector<pair<int,int>> edges;
    while ((int)edges.size() < m) {
        // 变量：u=随机简单图候选边的第一个端点 u；v=随机简单图候选边的第二个端点 v。
        int u = randint(1, n), v = randint(1, n);
        if (u == v) continue;
        if (u > v) swap(u, v);
        if (used.count({u, v})) continue;
        used.insert({u, v});
        edges.push_back({u, v});
    }
    shuffle(edges.begin(), edges.end(), rng);
    for (auto& [u, v] : edges)
        if (randint(0, 1)) swap(u, v);
    return edges;
}

// 生成 n 个不重复的随机整数（值域 [lo, hi]）
// 接口：rand_distinct(n, lo, hi)：生成范围内互不相同的随机整数；返回类型 vector<int>。
// 参数：需要生成的整数个数 n；随机取值闭区间下界 lo；随机取值闭区间上界 hi。
vector<int> rand_distinct(int n, int lo, int hi) {
    // 变量：used=记录每个对象是否已被使用的标记 used。
    set<int> used;
    // 变量：res=尚未重复采样得到的随机整数结果 res。
    vector<int> res;
    while ((int)res.size() < n) {
        // 变量：x=当前采样并检查去重的随机整数 x。
        int x = randint(lo, hi);
        if (!used.count(x)) {
            used.insert(x);
            res.push_back(x);
        }
    }
    return res;
}

// 生成随机排列 [1, 2, ..., n]
// 接口：rand_perm(n)：生成 1..n 的随机排列；返回类型 vector<int>。
// 参数：随机排列的长度 n。
vector<int> rand_perm(int n) {
    // 变量：p=待随机打乱为 1..n 排列的数组 p。
    vector<int> p(n);
    iota(p.begin(), p.end(), 1);
    shuffle(p.begin(), p.end(), rng);
    return p;
}

// 接口：main(argc, argv)：用 argv[1] 初始化随机种子，生成并输出一组小规模随机数组；入口函数；每组测试前按注释清空全局状态。
// 参数：命令行参数个数 argc；命令行参数数组 argv。
int main(int argc, char* argv[]) {
    // 用命令行参数做种子
    rng.seed(atoi(argv[1]));

    // ===== 根据题目修改以下内容 =====
    int n = randint(1, 10); // 对拍时用小范围
    cout << n << "\n";
    // 变量：i=随机数据生成器当前输出的数组位置 i。
    for (int i = 0; i < n; i++) {
        cout << randint(1, 100);
        if (i + 1 < n) cout << " ";
    }
    cout << "\n";

    return 0;
}
\end{lstlisting}

\begin{tip}
\begin{itemize}
  \item \lstinline|mt19937| 比 \lstinline|rand()| 质量更好且可移植。用 \lstinline|mt19937_64| 生成 64 位随机数。
  \item \lstinline|uniform_int_distribution<int>(lo, hi)| 生成 $[lo, hi]$ 闭区间均匀分布。
  \item 生成随机树时，如果需要链状树（极端数据），让每个节点连接 $i-1$；如果需要菊花图，让所有节点连接节点 1。
  \item 随机图的边数 $m$ 不能超过 $\frac{n(n-1)}{2}$，否则死循环。
\end{itemize}
\end{tip}

\section{暴力解法（brute.cpp）}

\begin{note}
暴力程序的原则：正确性第一，效率无所谓。能用 $O(n^3)$ 就用 $O(n^3)$，能枚举就枚举，绝不优化。
\end{note}

\begin{lstlisting}
#include <bits/stdc++.h>
using namespace std;
#define int long long
#define endl '\n'

// 接口：solve()：读取一个整数数组，暴力枚举全部子数组并输出满足题目条件的计数；无返回值。
void solve() {
    // 变量：n=输入数组的元素个数 n。
    int n;
    cin >> n;
    // 变量：a=暴力解法读入的长度为 n 的整数数组 a。
    vector<int> a(n);
    // 变量：i=当前从输入读取数组元素的下标 i。
    for (int i = 0; i < n; i++) cin >> a[i];

    // ===== 暴力做法（保证正确，不管复杂度）=====
    // 例：暴力枚举所有子集 / 全排列 / O(n^2) 枚举
    // 变量：ans=暴力枚举中累计满足题目条件的方案数 ans。
    int ans = 0;
    // 变量：i=暴力枚举子数组的起点 i。
    for (int i = 0; i < n; i++)
        // 变量：j=暴力枚举中从起点 i 向右扩展的子数组终点 j。
        for (int j = i; j < n; j++) {
            // ...
        }

    cout << ans << endl;
}

// 接口：main()：开启快速 I/O，并按测试数 t 调用暴力解 solve；入口函数；每组测试前按注释清空全局状态。
signed main() {
    ios::sync_with_stdio(0);
    cin.tie(0);
    // 变量：t=当前程序要执行的测试用例数量 t。
    int t = 1;
    // cin >> t; // 根据题目决定是否多测
    while (t--) solve();
    return 0;
}
\end{lstlisting}

\section{交互题对拍}

\begin{note}
交互题无法用管道对拍，需要写一个 interactor 同时充当评测端和数据生成器。
\end{note}

\begin{lstlisting}
// interactor.cpp —— 模拟交互评测
// 编译后用管道连接：
// Linux:  mkfifo pipe; (./interactor < pipe | ./sol > pipe)
// 或者直接写成单文件 interactor，用 popen 启动 sol
#include <bits/stdc++.h>
using namespace std;

// 接口：main()：生成隐藏答案，读取选手猜测并回复比较结果，直至猜中或超过询问上限；入口函数；每组测试前按注释清空全局状态。
int main() {
    mt19937 rng(42);
    // 生成隐藏数据
    // 变量：secret=交互题中仅供本地模拟的隐藏答案 secret。
    int secret = uniform_int_distribution<int>(1, 1000)(rng);

    // 模拟交互过程
    // 变量：queries=输入询问集合 queries。
    int queries = 0;
    while (true) {
        // 变量：guess=本轮向交互器提交的猜测值 guess。
        int guess;
        // 读取选手的输出（选手的 cout 是这里的 cin）
        if (!(cin >> guess)) break;
        queries++;
        if (guess == secret) {
            cout << "=" << endl; // 正确
            break;
        } else if (guess < secret) {
            cout << "<" << endl; // 太小
        } else {
            cout << ">" << endl; // 太大
        }
    }
    // 检查询问次数是否超限
    cerr << "Queries used: " << queries << endl;
    return 0;
}
\end{lstlisting}

\begin{warn}
\begin{itemize}
  \item 对拍时 gen 的数据范围要小！暴力程序通常是 $O(n^2)$ 或 $O(2^n)$，$n$ 太大跑不完。一般 $n \le 10$ 起步，跑上万组。
  \item 每次对拍脚本传入不同的种子 \lstinline|./gen $i|，确保数据多样性。找到 WA 后用同一个种子可以复现。
  \item 注意行末空格和换行的差异可能导致 diff 报 WA，可以在对比前做 trim。
  \item 如果正解涉及浮点数输出，diff 无法判断精度误差，需要写 checker 比较相对/绝对误差。
  \item 生成的输入必须严格满足题目约束（如连通图、无重边、值域范围等），否则暴力和正解可能行为都不确定。
\end{itemize}
\end{warn}

% ==========================================================
\chapter{竞赛环境速查}
% ==========================================================

\section{Windows 命令行（CMD / PowerShell）}

\subsection{编译与运行}

\begin{lstlisting}
:: 基础编译
g++ -o sol.exe sol.cpp

:: 推荐编译选项（竞赛常用）
g++ -O2 -std=c++17 -o sol.exe sol.cpp

:: 开所有警告 + 调试信息（调试时用）
g++ -Wall -Wextra -Wshadow -g -fsanitize=address,undefined -o sol.exe sol.cpp

:: 编译并立即运行
g++ -O2 -o sol.exe sol.cpp && sol.exe

:: 从文件读入、输出到文件
sol.exe < in.txt
sol.exe < in.txt > out.txt

:: 同时输出到屏幕和文件（PowerShell）
.\sol.exe < in.txt | Tee-Object out.txt

:: 计时运行（PowerShell）
Measure-Command { .\sol.exe < in.txt > out.txt }
\end{lstlisting}

\begin{warn}
\begin{itemize}
  \item \lstinline|-fsanitize=address,undefined| 可以检测数组越界、未定义行为等，调试利器！但会显著变慢，提交前务必去掉。
  \item \lstinline|-O2| 和 \lstinline|-fsanitize| 不要同时用，优化可能把 sanitizer 的检测优化掉。
  \item Windows 下 PowerShell 和 CMD 的重定向语法一样，但管道行为有差异。推荐用 PowerShell。
\end{itemize}
\end{warn}

\subsection{文件操作}

\begin{lstlisting}
:: 查看文件内容
type in.txt                  :: CMD
cat in.txt                   :: PowerShell (Get-Content 的别名)

:: 比较两个文件
fc out_brute.txt out_sol.txt :: CMD（按行比较）
diff out1.txt out2.txt       :: PowerShell (Compare-Object 的别名，需 Git)

:: 复制 / 移动 / 删除
copy a.txt b.txt
move a.txt b.txt
del a.txt

:: 创建目录
mkdir contest

:: 列出当前目录文件
dir                          :: CMD
ls                           :: PowerShell
\end{lstlisting}

\subsection{进程与调试}

\begin{lstlisting}
:: 强制终止程序（死循环时用）
Ctrl + C                     :: 终端中直接按

:: 查看正在运行的进程
tasklist | findstr sol.exe

:: 强制杀死进程
taskkill /F /IM sol.exe

:: 查看程序返回值（上一条命令的退出码）
echo %errorlevel%            :: CMD
echo $LASTEXITCODE           :: PowerShell
\end{lstlisting}

\subsection{常用技巧}

\begin{lstlisting}
:: 快速生成测试数据并对拍（单行命令）
g++ -O2 -o gen.exe gen.cpp && g++ -O2 -o brute.exe brute.cpp && g++ -O2 -o sol.exe sol.cpp

:: 批量编译当前目录所有 cpp 文件（PowerShell）
Get-ChildItem *.cpp | ForEach-Object {
    g++ -O2 -o ($_.BaseName + ".exe") $_.FullName
}

:: 清理编译产物（PowerShell）
Remove-Item *.exe, *.o, in.txt, out*.txt -ErrorAction SilentlyContinue

:: 查看 g++ 版本
g++ --version

:: 检查是否安装了某工具
where g++
where python
\end{lstlisting}

\begin{tip}
\begin{itemize}
  \item 文件重定向 \lstinline|< in.txt > out.txt| 在比赛中比在代码里写 \lstinline|freopen| 更灵活，不用改代码。
  \item 如果需要在代码内重定向（某些 OJ 要求），用：\\
    \lstinline|freopen("input.txt", "r", stdin);|\\
    \lstinline|freopen("output.txt", "w", stdout);|
  \item PowerShell 中 \lstinline|Measure-Command| 可以精确测量运行时间，判断是否会 TLE。
\end{itemize}
\end{tip}

\section{VS Code 快捷键速查}

\begin{note}
以下快捷键基于 Windows 默认键位。macOS 把 \lstinline|Ctrl| 换成 \lstinline|Cmd|。
\end{note}

\subsection{文件与编辑}

\begin{lstlisting}
Ctrl + N              新建文件
Ctrl + O              打开文件
Ctrl + S              保存
Ctrl + Shift + S      另存为
Ctrl + W              关闭当前标签
Ctrl + Tab            切换已打开的文件

Ctrl + Z              撤销
Ctrl + Shift + Z      重做（或 Ctrl + Y）
Ctrl + X              剪切整行（光标在行上，不选中也可以）
Ctrl + C              复制整行（同上）
Ctrl + Shift + K      删除整行
Alt + Up/Down         上下移动当前行
Shift + Alt + Up/Down 向上/下复制当前行
Ctrl + D              选中当前单词，再按选下一个相同的
Ctrl + Shift + L      选中所有相同的单词（批量重命名变量超好用）
\end{lstlisting}

\subsection{查找与替换}

\begin{lstlisting}
Ctrl + F              查找
Ctrl + H              查找并替换
Ctrl + Shift + F      全局搜索（跨文件）
Ctrl + G              跳转到指定行号
F12                   跳转到定义
Alt + F12             速览定义（弹窗）
Ctrl + Shift + O      跳转到当前文件的符号（函数名等）
\end{lstlisting}

\subsection{多光标与选择}

\begin{lstlisting}
Alt + Click           添加多光标
Ctrl + Alt + Up/Down  上下添加光标（竖向多行编辑）
Ctrl + Shift + L      选中所有匹配项并添加光标
Ctrl + L              选中当前行
Shift + Alt + 拖拽    列选择（矩形选择）
Ctrl + Shift + [/]    折叠/展开代码块
\end{lstlisting}

\subsection{终端}

\begin{lstlisting}
Ctrl + `              打开/关闭集成终端
Ctrl + Shift + `      新建终端
Ctrl + Shift + 5      拆分终端（左右并排）
Ctrl + PageUp/Down    切换终端标签
\end{lstlisting}

\begin{tip}
\begin{itemize}
  \item 在终端中按 \lstinline|Up| 键可以翻出上一条命令，快速重复编译运行。
  \item 可以配置 VS Code 的 \lstinline|tasks.json| 实现一键编译运行（见下方配置）。
\end{itemize}
\end{tip}

\subsection{调试}

\begin{lstlisting}
F5                    启动调试
Shift + F5            停止调试
F9                    切换断点
F10                   单步跳过（Step Over）
F11                   单步进入（Step Into）
Shift + F11           单步跳出（Step Out）
Ctrl + Shift + D      打开调试面板
\end{lstlisting}

\subsection{其他实用操作}

\begin{lstlisting}
Ctrl + Shift + P      命令面板（万能入口）
Ctrl + ,              打开设置
Ctrl + B              切换侧边栏显示/隐藏
Ctrl + J              切换底部面板显示/隐藏
Ctrl + K, Ctrl + S    查看所有快捷键
Ctrl + /              注释/取消注释当前行
Ctrl + Shift + A      块注释（/* ... */）
Ctrl + K, Ctrl + F    格式化选中代码
Shift + Alt + F       格式化整个文件
Ctrl + Shift + V      预览 Markdown
\end{lstlisting}

\section{VS Code 竞赛配置}

\subsection{一键编译运行（tasks.json）}

\begin{note}
在 \lstinline|.vscode/tasks.json| 中配置，按 \lstinline|Ctrl+Shift+B| 即可一键编译运行当前文件。
\end{note}

\begin{lstlisting}
{
    "version": "2.0.0",
    "tasks": [
        {
            "label": "compile and run",
            "type": "shell",
            "command": "g++",
            "args": [
                "-O2", "-std=c++17",
                "-Wall", "-Wextra",
                "-o", "${fileDirname}/${fileBasenameNoExtension}.exe",
                "${file}"
            ],
            "group": {
                "kind": "build",
                "isDefault": true
            },
            "problemMatcher": "$gcc"
        }
    ]
}
\end{lstlisting}

\subsection{调试配置（launch.json）}

\begin{lstlisting}
{
    "version": "0.2.0",
    "configurations": [
        {
            "name": "Debug C++",
            "type": "cppdbg",
            "request": "launch",
            "program": "${fileDirname}/${fileBasenameNoExtension}.exe",
            "args": [],
            "cwd": "${fileDirname}",
            "stopAtEntry": false,
            "externalConsole": true,
            "MIMode": "gdb",
            "miDebuggerPath": "gdb",
            "preLaunchTask": "compile and run"
        }
    ]
}
\end{lstlisting}

\subsection{推荐插件}

\begin{note}
\begin{itemize}
  \item \textbf{C/C++}（Microsoft）：语法高亮、智能补全、调试支持。
  \item \textbf{Code Runner}：右键一键运行代码，支持输入重定向。
  \item \textbf{Competitive Programming Helper (cph)}：自动抓取样例、一键测试，支持 Codeforces/AtCoder 等。
  \item \textbf{Error Lens}：将编译错误/警告直接显示在代码行内，快速定位问题。
  \item \textbf{Bracket Pair Colorizer}（VS Code 已内置）：括号彩色匹配，嵌套多层时一目了然。
\end{itemize}
\end{note}

\subsection{推荐用户设置（settings.json 片段）}

\begin{lstlisting}
{
    // 自动保存
    "files.autoSave": "afterDelay",
    "files.autoSaveDelay": 1000,
    // 字体推荐等宽字体
    "editor.fontFamily": "Consolas, 'Courier New', monospace",
    "editor.fontSize": 15,
    // 显示行号
    "editor.lineNumbers": "on",
    // Tab 转空格，宽度 4
    "editor.tabSize": 4,
    "editor.insertSpaces": true,
    // 括号匹配高亮
    "editor.bracketPairColorization.enabled": true,
    // 终端默认 PowerShell
    "terminal.integrated.defaultProfile.windows": "PowerShell"
}
\end{lstlisting}

\begin{tip}
\begin{itemize}
  \item 比赛时建议关闭自动补全的延迟弹出，避免打字时频繁弹出提示干扰思路：\\
    \lstinline|"editor.quickSuggestions": false|
  \item 使用 \lstinline|Ctrl+Shift+P| 输入 ``Snippets'' 可以自定义代码片段，把常用模板设为一键插入。
  \item 如果终端太小，\lstinline|Ctrl+`| 打开终端后，拖拽顶部边框可以调整高度。
\end{itemize}
\end{tip}

\inputnewboard{remake/large/12_索引.tex}

\end{document}
